% Generated draft: requires mathematical and source review.
\documentclass[11pt]{article}
\usepackage[margin=25mm]{geometry}
\usepackage{fontspec}
\setmainfont{Latin Modern Roman}
\usepackage{amsmath,amssymb,mathtools,unicode-math}
\setmathfont{Latin Modern Math}
\usepackage{enumitem,needspace,xurl,hyperref}
\setlength{\parindent}{0pt}
\setlength{\parskip}{5pt}
\setlength{\emergencystretch}{3em}
\newcommand{\N}{\mathbb{N}}
\newcommand{\Z}{\mathbb{Z}}
\newcommand{\Q}{\mathbb{Q}}
\newcommand{\R}{\mathbb{R}}
\newcommand{\C}{\mathbb{C}}
\newcommand{\F}{\mathbb{F}}
\newcommand{\A}{\mathbb{A}}
\newcommand{\PP}{\mathbb{P}}
\newcommand{\E}{\mathbb{E}}
\newcommand{\eps}{\varepsilon}
\newcommand{\dd}{\,\mathrm{d}}
\begin{document}
\section*{UnsolvedMath ID: 801049}
\textbf{Source problem code:} catalogue entry
\subsection{Definitions and mathematical statement}
The supplied Question 1 asks whether a H\'enon map of some degree admits a wild attractor, described there as a Cantor set attracting a set of positive Lebesgue measure and strictly contained in a transitive set. The excerpt does not specify the ambient field, the class of H\'enon maps, or the convention for transitivity, so it does not determine a unique fully specified existence statement.

A conditional formulation is as follows. Fix an ambient Euclidean phase space \(X\), its metric \(d\) and Lebesgue measure \(m_X\), and an intended class \(\mathcal H(X)\) of H\'enon maps. A Cantor set means a nonempty compact, perfect, totally disconnected subset of \(X\). For a map \(F\colon X\to X\) and a Cantor set \(K\subseteq X\), define its basin of attraction by
\[
 B_F(K)=\left\{z\in X:\lim_{n\to\infty}\operatorname{dist}(F^n(z),K)=0\right\},
 \qquad \operatorname{dist}(z,K)=\inf_{w\in K}d(z,w).
\]
Under the convention that an attractor is invariant and a transitive set is an invariant set on which the restriction is topologically transitive, the question becomes: do there exist \(F\in\mathcal H(X)\), of some degree, and sets \(K,T\subseteq X\) such that
\[
\begin{gathered}
 K\text{ is a Cantor set},\qquad F(K)=K,\qquad m_X(B_F(K))>0,\\
 K\subsetneq T,\qquad F(T)=T,
\end{gathered}
\]
and, for every pair of nonempty relatively open subsets \(U,V\subseteq T\), there exists an integer \(n\geq 0\) such that
\[
 F^n(U)\cap V\neq\varnothing?
\]
Here the choices of \(X\) and \(\mathcal H(X)\) remain unresolved. The basin, invariance, and topological-transitivity conventions above make one reading explicit; they require confirmation against the complete source context. The excerpt also does not state whether \(T\) must be closed or compact, or whether further conditions on the attractor or its basin are intended.
\subsection{Short English statement}
Berger and Bedford ask whether some H\'enon map, with degree unrestricted in the excerpt, has an invariant Cantor set attracting a set of positive Lebesgue measure while lying strictly inside a transitive invariant set. The ambient space, admissible map class, and precise dynamical conventions require source clarification. The registry labels the problem open; this is a supplied status, not an independently verified assessment.
\subsection{Sources}
\begin{itemize}
\item Local UnsolvedMath registry, permanent ID \(801049\), source problem code \texttt{catalogue entry}, title ``Henon Maps: Open Problems --- Question 1.'' The supplied record attributes the problem to Berger and Bedford, with proposed year 2023, and labels its status \texttt{open}.
\item Berger, Bedford, \emph{H\'enon Maps: A List of Open Problems} (2023), Question 1. Supplied source URL: \url{https://arxiv.org/abs/2312.03907}. The background records access on 2026-10-16 and extraction by \texttt{arxiv\_latex}. These are supplied provenance details; the source and its current status have not been independently verified here.
\end{itemize}
\end{document}
